\chapter{Conclusiones y líneas futuras}
\label{ch:conclusiones}

Este trabajo partió de una observación sencilla: la elección de tutor académico es un problema de recuperación de información que las universidades tienen sin resolver pese a publicar los datos necesarios para resolverlo.
Sobre esa base se ha diseñado, construido y evaluado \proyecto, un sistema completo que convierte el Portal Científico de la \gls{upm} en un servicio de recomendación de tutores consultable en lenguaje natural.

Retomando los objetivos formulados en la Sección~\ref{sec:objetivos}:

\begin{itemize}
    \item \textbf{\ref{obj:adquisicion} - Adquisición y estructuración.}
    Se ha construido un pipeline automatizado e incremental que censa, descarga y estructura la actividad pública de la \gls{upm} en un grafo de investigadores, publicaciones, proyectos y supervisiones, con la fuente cruda conservada para reprocesado.
    El corpus resultante reúne 172\,621 nodos y 268\,886 relaciones, con 709 investigadores efectivamente indexados de los 996 que componen el perímetro de las cuatro escuelas seleccionadas para las búsquedas.
    \item \textbf{\ref{obj:perfilado} - Perfilado semántico.}
    Cada investigador dispone de un perfil enriquecido (temas normalizados, dominios de especialización y biografía investigadora generada incrementalmente con un \gls{llm} local) que agrega toda su evidencia pública en una representación homogénea.
    \item \textbf{\ref{obj:emparejamiento} - Emparejamiento.}
    El sistema implementa cuatro estrategias de recuperación operativas (léxica con \gls{llm}, semántica directa, semántica enriquecida e híbrida con fusión \gls{rrf}) que atienden consultas libres en cualquier idioma.
    \item \textbf{\ref{obj:evaluacion} - Evaluación objetiva.}
    Se ha definido y aplicado un marco de evaluación reproducible que usa la supervisión histórica como verdad de referencia, eliminando el etiquetado manual.
    La comparación muestra que la estrategia híbrida (fusión \gls{rrf} de la señal dispersa y la densa) es la mejor en todas las métricas (\gls{mrr} 0,153), por delante de la semántica enriquecida (0,146) y de la directa (0,137), y que la recuperación semántica supera con creces a la léxica (la mejor estrategia léxica alcanza 0,091, frente a 0,153 de la híbrida).
    El modelo de lenguaje en consulta aporta en ambas vías.
    \item \textbf{\ref{obj:prototipo} - Prototipo.}
    El prototipo web es funcional, explica sus recomendaciones y se despliega íntegramente \emph{on-premise} con componentes abiertos, sin dependencia de servicios externos ni coste por consulta.
\end{itemize}

En conjunto, la respuesta a la pregunta de partida es afirmativa con matices.
El dato público institucional, procesado íntegramente con \glspl{llm} en infraestructura propia, basta para construir un servicio de recomendación de tutores que es a la vez \emph{útil} (la fusión híbrida sitúa al supervisor real entre los primeros puestos con una frecuencia muy superior a la de una búsqueda léxica ingenua) y \emph{medible}, gracias a un marco de evaluación reproducible que no precisa etiquetado manual.
Los matices son dos: las cifras absolutas son modestas y deben leerse como cotas inferiores de la utilidad real (Sección~\ref{sec:eval-validez}) y la viabilidad descansa sobre un coste de cómputo no despreciable (Sección~\ref{sec:disc-coste-llm})
Ninguno invalida la tesis de fondo: que la recomendación de tutores puede abordarse hoy, de forma soberana y transparente, con tecnología abierta y datos que la institución ya publica.

\section{Contribuciones}
\label{sec:conc-contribuciones}

Más allá del sistema concreto, el trabajo deja cuatro aportaciones reutilizables: (i) un corpus estructurado y reconstruible de la actividad investigadora de las escuelas de la \gls{upm}; (ii) una metodología de evaluación sin etiquetado manual aplicable a cualquier institución con registros de supervisión; (iii) una comparación empírica de estrategias de recuperación representativas del estado del arte sobre un corpus institucional real; y (iv) un prototipo de código abierto que demuestra la viabilidad del despliegue local extremo a extremo.

\section{Líneas futuras}
\label{sec:conc-futuro}

El trabajo abre varias líneas naturales de continuación, ordenadas aquí de la más inmediata a la más ambiciosa:

\begin{mydescription}
    \item[Detección de colaboraciones entre investigadores.]
    El grafo y los perfiles semánticos ya construidos permiten plantear la pregunta simétrica a la de este \gls{tfm}: qué pares de investigadores presentan afinidad temática sin historial de colaboración.
    Es la extensión prevista desde el diseño (Sección~\ref{sec:objetivos}) y uno de los objetivos institucionales que enmarcan el trabajo: fomentar colaboraciones interdepartamentales mediante el análisis de las redes de coautoría (Sección~\ref{sec:contexto}).
    \item[Evaluación con usuarios reales.]
    Complementar la evaluación con la prueba piloto con estudiantes ya prevista como trabajo futuro (Sección~\ref{sec:contexto}): consultas genuinas y la valoración directa de las recomendaciones recibidas.
    Esto permitiría, además, cuantificar el sesgo del marco planteado (Sección~\ref{sec:eval-validez}).
    \item[Mejoras de recuperación.]
    Dos candidatas claras: el \emph{re-ranking} de los primeros resultados con un modelo más potente (\emph{cross-encoder} o \gls{llm} como juez) y el ajuste fino del modelo de embedding al dominio académico bilingüe utilizando los propios pares $supervisión-proyecto$ como señal de entrenamiento.
    \item[Equidad y arranque en frío.]
    Incorporar mecanismos que den visibilidad al profesorado con poco historial (desde la exploración controlada en el ranking hasta perfiles basados en su producción doctoral) y métricas de equidad en el banco de pruebas.
    \item[Ampliación del perímetro.]
    Extender el corpus al conjunto de la \gls{upm} y validar la portabilidad del flujo en otra institución, idealmente con un portal de estructura distinta, para convertir la generalización argumentada (Sección~\ref{sec:disc-generalizacion}) en generalización demostrada.
    Esta ampliación es ante todo un ejercicio de planificación de capacidad: dado el coste de generación medido (Sección~\ref{sec:disc-coste-llm}), procesar toda la universidad exige dimensionar de antemano el cómputo necesario, ya sea ampliando los recursos propios o recurriendo a un servicio de mayor potencia (esta última opción en tensión con el requisito de soberanía de \ref{req:soberania}).
    Con esa previsión, el salto de escala es manejable y no un obstáculo de fondo.
    \item[Integración institucional.]
    Llevar el prototipo a un servicio: autenticación corporativa, retroalimentación de los usuarios como señal de mejora continua y conexión con los procesos reales de oferta y asignación de trabajos.
    \item[Base de conocimiento conversacional y agentes.]
    Los perfiles semánticos construidos en este trabajo son legibles tanto por algoritmos de recuperación como por modelos generativos (Sección~\ref{sec:sota-posicionamiento}), lo que abre la línea más ambiciosa: exponer la actividad investigadora de la institución como una base de conocimiento consultable por \glspl{llm}.
    La vía técnica natural es un servidor \gls{mcp} que publique los perfiles y las operaciones de búsqueda como herramientas estandarizadas, de modo que un asistente conversacional (o un sistema multiagente) pueda razonar sobre el corpus, responder preguntas en lenguaje natural y orquestar tareas académicas (sugerir tribunales, detectar sinergias, redactar justificaciones) sin acceso directo a las bases de datos subyacentes.
\end{mydescription}
